% `tex.web` §336-§339: what happens when a balanced scan meets something that
% may not be inside it.
%
%     @<Tell the user what has run away and try to recover@>=
%     begin runaway;
%     if cur_cs=0 then print_err("File ended")
%     else begin cur_cs:=0; print_err("Forbidden control sequence found"); end;
%     print(" while scanning "); … print(" of "); sprint_cs(warning_index);
%
% §305's `scanner_status` says what is being scanned and so what the message
% calls it — a `definition`, a `use`, a `preamble`, a `text` — and
% `warning_index` says whose. §306 prints what had been absorbed so far under a
% `Runaway text?` heading, §337 backs the offending name up so it can be reread
% *as a command*, and §339 inserts the `}` that was missing.
%
% The recovery has one more step that is easy to miss: §337 ends by turning the
% offending token into a space, and the scan carries on with it. So the token
% list here ends `a<space>` and not `a`, the inserted `}` closes the group on
% the next turn of the loop, and the `\outer` name arrives afterwards as an
% ordinary command. Breaking out of the scan instead leaves that `}` to
% `main_control`, which meets it with nothing open and says `Too many }'s` a
% second time.
%
% Until this existed, `\outer` was recorded and never enforced, and a runaway
% swallowed the rest of the file. trip.tex line 354 does exactly that on
% purpose — `\tokens{\a^^@^^@a\par!` with an `\outer` name inside and no
% closing brace — so everything after it was never read: ninety lines of
% trip.tex and eight of its fourteen shipped pages.
%
% §262's `print_cs` prints the name through §48's string pool, in which every
% character below 32 *is* the two-character string `^^X`. So a name made of
% control characters — trip.tex's `\a^^@^^@a`, whose middle two characters are
% NUL — is quoted readably rather than written raw into the log. No case here
% builds such a name: reproducing one needs NUL to be a letter, and getting
% *that* right is its own divergence (ours keeps the NULs in the name where
% TeX does not). trip.log is where this half was measured.
\catcode`\{=1 \catcode`\}=2 \catcode`\#=6
\outer\def\bad{X}
\toksdef\tokens=9

% Absorbing a token list: `text of \tokens`.
\tokens{a\bad b}
\message{PROBE:one=[\the\tokens]}

% A macro body: `definition of \worse`.
\def\worse{p\bad q}
\message{PROBE:two=[\meaning\worse]}

% ...and an ordinary use of the same name is fine.
\bad
\message{PROBE:three=ok}
\end
